\section{Limitations}

The synthetic absent benchmark is useful for controlled evaluation, but it is not a substitute for a larger human-authored absent-query dataset. The 100-example realistic validation set is an external-validity check, not a final benchmark.

The OCR component is brittle for small, stylized, low-contrast, and non-text UI elements. The current non-text/image-cue results should therefore be treated as a proxy experiment rather than evidence of full non-text target support.

Evidence-aware VLM verification is promising and now has matched image-split validation for SeekUI, but the real-absent evidence-aware variant still needs a direct validation run with per-example evidence and predictions. The same evidence-aware check should also be run on SeekUI-SFT.

The paper should therefore separate the interpretable-method claim from the best-verifier claim. Combined AND remains the split-controlled interpretable method, while evidence-aware VLM is the strongest validated verifier currently available.

Finally, the current approach is post-hoc. It changes the present/absent decision after SeekUI has generated a scanpath; it does not train a model to produce better target-absent search behavior. A stronger future model should jointly generate search behavior, uncertainty, and stopping decisions.
